\clearpage
\VerdierPhysicalPageAnchor{0194}%
\begingroup\small\setlength{\parskip}{0pt}\setlength{\emergencystretch}{1.5em}
\setlength{\abovedisplayskip}{4pt plus 2pt minus 2pt}%
\setlength{\belowdisplayskip}{4pt plus 2pt minus 2pt}%

\VerdierSemanticAnchor{ch3d-2-2-6-proof-cone-cont}
($L_0$ is the component of degree $0$ of $c(m)$). The cohomology of
$c(m)$ vanishes in every degree other than $-1$, and equals $X$ in degree
$-1$. The morphism $L_n\longrightarrow L_n\oplus Z_{n-1}$ whose matrix
has the form $\begin{pmatrix}\operatorname{id}\\ *\end{pmatrix}$ admits a
retraction. There is therefore a morphism
$v\colon Z_{n-1}\longrightarrow L_{n-1}\oplus L'_{n-2}$ such that the
complex

\VerdierSemanticAnchor{ch3d-2-2-6-reduced-cone}
\[
\begin{aligned}
\cdots\longrightarrow0\longrightarrow0\longrightarrow Z_{n-1}
\xrightarrow{v}L_{n-1}\oplus L'_{n-2}\longrightarrow\cdots
\longrightarrow L_1\oplus L'_0\\
\longrightarrow L_0\longrightarrow0\longrightarrow\cdots
\end{aligned}
\]
is quasi-isomorphic to $c(m)$. On the other hand, by property $(E_4)$ the
kernel of $L_1\oplus L'_0\longrightarrow L_0$ belongs to $M$. Denote this
kernel by $L''_1$. The complex

\VerdierSemanticAnchor{ch3d-2-2-6-resolution-x}
\[
\begin{aligned}
0\longrightarrow0\longrightarrow Z_{n-1}
\longrightarrow L_{n-1}\oplus L'_{n-2}\longrightarrow\cdots
\longrightarrow L_2\oplus L'_1\\
\longrightarrow L''_1\longrightarrow X\longrightarrow0\longrightarrow\cdots
\end{aligned}
\]
is therefore acyclic. We have thus constructed a resolution of $X$ of
length $n-1$, and consequently $M\text{-}\dim(X)\leq n-1$. On the other
hand, it is clear that
$M\text{-}\dim(X)\geq M\text{-}\dim(Y)-1$.

\VerdierSemanticAnchor{ch3d-2-2-6-conclusion}
This completes the proof of Lemma (2.2.6) and Proposition (2.2.7).

\VerdierSemanticAnchor{ch3d-2-3}
\medskip\noindent\textbf{2.3. Some lemmas on resolutions.}

\VerdierSemanticAnchor{ch3d-2-3-1}
\medskip\noindent\textbf{Proposition 2.3.1.} \textbf{[22]}\quad
\textit{a) Let $M$ be a set of objects of $\mathcal A$ satisfying
properties $(E_1)$, $(E_2)$, and $(E_3)$ (2.2.1). Let
$Y\in\operatorname{Ob}\operatorname{comp}^{-}(\mathcal A)$ be a complex
whose components belong to $M^\wedge$ (2.2.1), and let
$m\colon X\longrightarrow Y$ be a quasi-isomorphism. There is a complex
$L^\bullet\in\operatorname{Ob}\operatorname{comp}^{-}(\mathcal A)$ whose
components belong to $M$, and a quasi-isomorphism
$w\colon L^\bullet\longrightarrow X$.}

\noindent\textit{b) Conversely, let $M$ be a set of objects of
$\mathcal A$ satisfying property $(E_1)$. Suppose that for every complex
$Y\in\operatorname{Ob}\operatorname{comp}^{-}(\mathcal A)$ whose
components belong to $M$, and every quasi-isomorphism
$m\colon X\longrightarrow Y$, there is a complex $L^\bullet$ whose
components belong to $M$ and a quasi-isomorphism
$w\colon L^\bullet\longrightarrow X$. Then $M$ satisfies properties
$(E_2)$ and $(E_3)$.}

\VerdierSemanticAnchor{ch3d-2-3-1-proof}
Let us prove a). First, there is a complex
$X'\in\operatorname{Ob}\operatorname{comp}^{-}(\mathcal A)$ and a
quasi-isomorphism $X'\longrightarrow X$ (1.2.1). We may therefore reduce
to the case where $X$ is an object of
$\operatorname{comp}^{-}(\mathcal A)$. By (2.1.4) and (2.1.5), there is a
sequence of morphisms

\VerdierSemanticAnchor{ch3d-2-3-1-xn-sequence}
\[
u_n\colon X_n\longrightarrow X_{n-1},
\qquad n\in\mathbf Z,
\]
\endgroup
